\section*{Chapitre \ref{ch:rc:build-a-risk-engine} --- Construction~: un moteur de risque}

\begin{solution}{exo:rc:build-a-risk-engine:1}
\omterm{def:rc:build-a-risk-engine:reval}{Réévaluation complète}~: 251 appels pour chacun des deux (la base et 250 scénarios). Grille~: $1+8\times6=49$ pour le swap (huit \omterm{def:rc:curve-construction:pillar}{piliers},
six points non nuls chacun) et $1+9\times6=55$ pour l'option (le spot inclus). Delta--gamma~: $1+2\times8=17$ et
$1+2\times9=19$. Portefeuille~: $1\,000\times251=251\,000$~; $800\times49+200\times55=50\,200$~; $800\times17+200\times19=17\,400$.
\end{solution}

\begin{solution}{exo:rc:build-a-risk-engine:2}
La grille valorise chaque transaction sur un ensemble fixe de chocs~; chaque scénario ne coûte alors qu'une interpolation et une somme,
et non un appel de valorisation. Ajouter des scénarios n'ajoute que de l'arithmétique (tant qu'ils restent dans le domaine de la grille).
\end{solution}

\begin{solution}{exo:rc:build-a-risk-engine:3}
La VaR est un quantile, non une somme~: la VaR d'un desk ne se calcule pas à partir des VaR de ses portefeuilles, ni celle de l'établissement à partir
de celles de ses desks. Conserver le P\&L de chaque transaction par scénario permet au moteur de calculer la mesure de n'importe quel nœud, et de tout nouveau
regroupement, en sommant d'abord les colonnes.
\end{solution}

\begin{solution}{exo:rc:build-a-risk-engine:4}
L'\omterm{def:rc:market-risk-measures:es}{expected shortfall} (ES) autonome des taux (47.34~millions de dollars) et du FX (14.63~millions de dollars) s'additionnent à 61.96~millions de dollars, sur les valeurs non arrondies, contre 50.89~millions de dollars pour l'établissement~: 11.07
millions de dollars de diversification. L'\omterm{def:rc:the-valuation-adjustment-desk:euler}{allocation d'Euler} attribue presque toute l'ES de l'établissement aux taux.
\end{solution}

\begin{solution}{exo:rc:build-a-risk-engine:5}
Le desk est vendeur d'options à environ un mois de leur échéance. Pour un mouvement de 2.8~\%, le gamma baisse à mesure que les options s'éloignent
de la monnaie, si bien que les pertes croissent plus lentement que la parabole qui extrapole le gamma du jour.
\end{solution}

\begin{solution}{exo:rc:build-a-risk-engine:6}
La valeur de la swaption dépend du taux swap à terme, somme pondérée des mouvements des \omterm{def:rc:curve-construction:pillar}{piliers}, et sa convexité réside dans
le carré de cette somme~: les produits croisés entre \omterm{def:rc:curve-construction:pillar}{piliers}. Les méthodes par facteur ne conservent que les carrés.
Remèdes peu coûteux~: une grille sur le taux swap à terme de la transaction (un facteur qui porte les termes croisés), une
grille de décalage parallèle plus deltas par \omterm{def:rc:curve-construction:pillar}{pilier} pour le reste, ou \omterm{def:rc:build-a-risk-engine:reval}{réévaluation complète} des quelques transactions concernées.
\end{solution}

\begin{solution}{exo:rc:build-a-risk-engine:7}
Options de taux en \omterm{def:rc:build-a-risk-engine:reval}{réévaluation complète}~: $200\times251 = 50\,200$~; options de change à longue échéance par delta--gamma~: $111\times19 = 2\,109$~; swaps
par grille~: $600\times49 = 29\,400$~; options de change à courte échéance par grille~: $89\times55 = 4\,895$. Total 86\,604, le nombre que le
compteur d'appels du moteur renvoie pour le plan.
\end{solution}

\begin{solution}{exo:rc:build-a-risk-engine:8}
L'erreur dépend du portefeuille et des scénarios~: de nouvelles options à courte échéance, un horizon plus long ou une période de stress
peut mettre en défaut une approximation qui avait passé le test. Le test doit être relancé chaque fois que l'un des deux change, et
régulièrement.
\end{solution}

\begin{solution}{pb:rc:build-a-risk-engine:1}
\textbf{1.} \omterm{def:rc:build-a-risk-engine:reval}{Réévaluation complète} 251\,000, grille 50\,200, delta--gamma 17\,400.
\textbf{2.} Appels $\times\,50\times0.02/16/3\,600$~: 4.36, 0.87 et 0.30 heures.
\textbf{3.} Grille et delta--gamma~; la \omterm{def:rc:build-a-risk-engine:reval}{réévaluation complète} manque une fenêtre de quatre heures.
\textbf{4.} La grille sur les transactions, les facteurs et les points~; la \omterm{def:rc:build-a-risk-engine:reval}{réévaluation complète} sur les transactions multipliées par les scénarios.
\textbf{5.} $4\times3\,600\times16/(50\,000\times0.02) = 230.4$ appels par transaction~: 229 scénarios.
\textbf{6.} 1.59~\% pour la grille, 0.74~\% pour le delta--gamma.
\textbf{7.} 8.64~\% pour la grille, 18.45~\% pour le delta--gamma.
\textbf{8.} Les options de change à courte échéance~: leur gamma varie sur un mouvement de dix jours.
\textbf{9.} Les swaptions~: leur convexité réside dans les produits croisés des mouvements des \omterm{def:rc:curve-construction:pillar}{piliers}.
\textbf{10.} Les mouvements croissent comme la racine carrée de l'horizon, et les termes négligés croissent plus vite que les
termes retenus.
\textbf{11.} Grille par défaut, \omterm{def:rc:build-a-risk-engine:reval}{réévaluation complète} pour les options de taux, delta--gamma pour les options de change à longue échéance.
\textbf{12.} 86\,604 appels, 1.50 heures.
\textbf{13.} 0.76~\%, sur le desk de change à courte échéance.
\textbf{14.} Oui~: son erreur delta--gamma est de 0.43~\%, car des options à plusieurs mois de leur échéance ont un gamma qui varie lentement~;
l'erreur de la grille y est plus grande (2.17~\%) car sept points interpolent mal un P\&L presque plat.
\textbf{15.} Les erreurs aux nœuds par rapport à la \omterm{def:rc:build-a-risk-engine:reval}{réévaluation complète}, selon un calendrier et chaque fois que le portefeuille ou les scénarios changent.
\textbf{16.} Le desk de change à courte échéance~: VaR à dix jours de 14.58~millions de dollars, pour une limite de 12~millions de dollars.
\textbf{17.} Que sa queue est une queue de taux~: les taux contribuent pour 47.34 sur les 50.89~millions de dollars d'ES.
\textbf{18.} Un calcul qui se termine à temps même en période de stress, avec des chiffres rapprochés et agrégés de façon largement
automatisée.
\textbf{19.} Résultat nommé~: \emph{l'échéance de six heures}~: \omterm{def:rc:build-a-risk-engine:reval}{réévaluation complète} 4.36 heures (exacte), grille 0.87 heures
(erreur maximale de 8.64~\%), delta--gamma 0.30 heures (18.45~\%), plan 1.50 heures (0.76~\%)~: à l'horizon de dix jours, le plan
est la méthode la moins coûteuse dans le budget de 2~\%~; à un jour, ce serait le delta--gamma.
\textbf{20.} Cela dépend du portefeuille, de l'horizon et des scénarios, qui changent tous.
\end{solution}

\begin{solution}{iq:rc:build-a-risk-engine:1}
Approvisionnement en transactions et en données de marché (instantanés, hiérarchie), génération des scénarios, réévaluation via la
bibliothèque de valorisation, tableau des P\&L, agrégation, mesures et limites, reporting et traçabilité~; orchestration des traitements par lots,
gestion des erreurs par transaction, et rapprochement avec les portefeuilles.
\iqlookfor{les étapes et le tableau des P\&L comme objet central.}
\end{solution}

\begin{solution}{iq:rc:build-a-risk-engine:2}
Mesurer où passe le temps~; paralléliser sur les transactions~; utiliser des grilles ou des sensibilités là où leur erreur est faible
et la \omterm{def:rc:build-a-risk-engine:reval}{réévaluation complète} là où elle ne l'est pas~; réutiliser les résultats des transactions inchangées~; mettre en cache les constructions de données de marché~; isoler chaque
échec de transaction sans arrêter le traitement.
\iqlookfor{mesure, parallélisme et un plan vérifié par rapport à la \omterm{def:rc:build-a-risk-engine:reval}{réévaluation complète}.}
\end{solution}

\begin{solution}{iq:rc:build-a-risk-engine:3}
La \omterm{def:rc:build-a-risk-engine:reval}{réévaluation complète} est exacte et coûte un appel par transaction et par scénario. Les grilles coûtent des appels par facteur et capturent
la non-linéarité à un seul facteur. Les sensibilités sont les moins coûteuses et échouent sur les grands mouvements et les effets croisés. Le choix
se mesure desk par desk.
\iqlookfor{coût, erreur et cas où chaque méthode échoue.}
\end{solution}

\begin{solution}{iq:rc:build-a-risk-engine:4}
La VaR est un quantile d'une somme, non une somme de quantiles, et elle n'est pas sous-additive en général. Allouer selon les
contributions d'Euler (la perte espérée de chaque desk dans les scénarios de queue de l'établissement pour l'ES), qui s'additionnent au total.
\iqlookfor{non-additivité et \omterm{def:rc:the-valuation-adjustment-desk:euler}{allocation d'Euler}.}
\end{solution}

\begin{solution}{iq:rc:build-a-risk-engine:5}
Instruments comparés aux formules fermées et aux limites~; approximations comparées à la \omterm{def:rc:build-a-risk-engine:reval}{réévaluation complète} sur de petits mouvements et sur
des scénarios réels~; agrégation sur des jeux de test et des identités (sommes d'Euler, totaux par nœud)~; jumeaux inter-langages sur
jeux de test partagés~; exécutions de non-régression sur sorties stockées~; tests de temps d'exécution.
\iqlookfor{une chaîne de tests, de la valorisation aux rapports.}
\end{solution}

\begin{solution}{iq:rc:build-a-risk-engine:6}
Gouvernance, architecture des données, données de risque exactes, complètes, disponibles à temps et adaptables, et rapports exacts, clairs,
exhaustifs. Les banques peinent en raison de systèmes hérités, de données fragmentées, de programmes sous-financés
et d'une attention insuffisante du conseil, comme l'ont relevé le rapport d'étape 2023 du Comité de Bâle et le guide 2024 de la BCE.
\iqlookfor{les principes et les obstacles pratiques.}
\end{solution}
